\chapter{Instrumentación Utilizada}
\label{chap:instrumentacion}

A lo largo del desarrollo de este trabajo práctico, se emplearon diversos instrumentos de medición y fuentes de excitación provistos en el laboratorio central. A continuación, se detallan las marcas, modelos y principales características de la instrumentación utilizada en los distintos experimentos:

\begin{itemize}
    \item \textbf{Osciloscopio de almacenamiento digital Rigol DS1052E}:
          Instrumento digital de 2 canales con un ancho de banda de $50\MHz$ y una velocidad de muestreo de hasta $1\text{ GS/s}$. Se utilizó para la visualización directa de las formas de onda de entrada y salida, la verificación de simetría en la máxima excursión (libre de distorsión) y la determinación del tiempo de crecimiento ($t_c$) mediante la función magnificadora horizontal ($X10$).
          
    \item \textbf{Generador de funciones GW Instek SFG-2120}:
          Generador de funciones sintetizado por DDS con un rango de frecuencia de hasta $20\MHz$. Proporcionó las señales de excitación senoidales y cuadradas de $1\kHz$ con tensión de salida controlada para inyectar a la entrada del amplificador.
          
    \item \textbf{Multímetro analógico UNIVO}:
          Instrumento de bobina móvil analógico provisto de escalas calibradas en decibeles ($\dBu$) referenciadas a una impedancia de $600\ohm$ (donde $0\dBu \equiv 0.775\V$). Se utilizó para medir los niveles absolutos de tensión en escala logarítmica para los cálculos de ganancia y potencia de salida.
          
    \item \textbf{Multímetro digital UNI-T UT33A+}:
          Instrumento digital auto-rango de $3\frac{3}{4}$ dígitos (resolución máxima de 3999 cuentas). Fue empleado para medir el valor de las resistencias de carga ($R_L$) y de resistores variables ($R_{e1}$, $R_{e2}$) para el cálculo de las impedancias de entrada y salida del circuito, así como para la posterior estimación de sus incertidumbres.
          
    \item \textbf{Fuente de alimentación de CC del Laboratorio Central}:
          Fuente de tensión regulada continua. Suministró la energía y polarización continua necesaria para el correcto funcionamiento de las etapas activas del amplificador bajo ensayo.
\end{itemize}
